\documentclass{apex_mdd}
\modelid{APEX-MDL-0003}
\mddversion{1.0}
\mddowner{Dr.\ Priya Nair (CRO)}
\mddvalidator{Dr.\ Samuel Achebe (Model Validation Officer)}
\mddstatus{Draft}
\reviewdate{2026-03-01}
\nextreview{2027-03-01}
\regulatory{BCBS d352 (Jan 2016), CRR3 2024, SR 11-7}

\title{FRTB Standardised Approach}
\begin{document}
\maketitle
\statusbadge

\tableofcontents
\newpage

\section{Model Overview}

\subsection{Purpose}
The FRTB Standardised Approach (SA) replaces the Basel 2.5 Standardised Approach with a sensitivity-based method (SBM) that better reflects risk. Under Basel IV / CRR3, all banks must report FRTB-SA capital; IMA banks use it as a floor. Apex uses FRTB-SA as the primary capital model pending IMA approval.

\subsection{Business Use}
\begin{itemize}
  \item \textbf{Regulatory capital (primary):} SA capital is the binding constraint until IMA approval.
  \item \textbf{Internal limit calibration:} SA sensitivities inform risk appetite for delta, vega, curvature.
  \item \textbf{Parallel run:} Compared against IMA output; material divergence triggers governance review.
\end{itemize}

\section{Theoretical Basis}

\subsection{Three Capital Components}

\subsubsection{1.\ Sensitivity-Based Method (SBM)}
Capital is computed from risk factor sensitivities across 7 risk classes (GIRR, CSR non-sec, CSR sec, FX, EQ, Commodity, VEGA).

\textbf{Delta capital:}
\begin{formulabox}
\[
  K_{\delta} = \sqrt{\sum_b \sum_{i,j\in b} \rho_{ij}\,WS_i\,WS_j}
\]
where $WS_i = RW_i \times s_i$ (risk weight $\times$ net sensitivity), $\rho_{ij}$ is the intra-bucket correlation.
\end{formulabox}

\textbf{Cross-bucket aggregation:}
\begin{formulabox}
\[
  K_{\mathrm{SBM}} = \sqrt{\sum_b K_b^2 + \sum_{b \neq c} \gamma_{bc}\,S_b\,S_c}
\]
where $\gamma_{bc}$ is the cross-bucket correlation and $S_b = \sum_i WS_i$ within bucket $b$.
\end{formulabox}

\subsubsection{2.\ Residual Risk Add-On (RRAO)}
\begin{formulabox}
\[
  \mathrm{RRAO} = 1.0\% \times N_{\mathrm{exotic}} + 0.1\% \times N_{\mathrm{other}}
\]
Applies to instruments with path-dependent payoffs, correlation products, and non-linear exotic optionality.
\end{formulabox}

\subsubsection{3.\ Default Risk Charge (DRC)}
\begin{formulabox}
\[
  \mathrm{DRC} = \max\!\Bigl(0,\;\sum_i \mathrm{JTD}_i \times \mathrm{LGD}_i \times RW_i\Bigr) \times W_{\mathrm{maturity}}
\]
Jump-to-default exposure for non-securitisation credit positions.
\end{formulabox}

\subsection{IR Tenor Buckets (GIRR)}
0.25Y, 0.5Y, 1Y, 2Y, 3Y, 5Y, 10Y, 15Y, 20Y, 30Y (10 buckets per currency).

\section{Mathematical Specification}

\subsection{Equity Risk Weights by Bucket}
\begin{center}
\small
\begin{tabular}{lll}
\toprule
\textbf{Bucket} & \textbf{Description} & \textbf{Delta RW} \\
\midrule
1  & Large-cap EM & 55\% \\
8  & Large-cap DM & 25\% \\
11 & Small-cap    & 70\% \\
\bottomrule
\end{tabular}
\end{center}

\subsection{IR Risk Weights (GIRR)}
Tenor-dependent: 1.7\% (3m) to 1.5\% (10Y), prescribed by BCBS d352 §21.52.

\textbf{Correlation parameters:} Three scenarios (high/medium/low) for conservative capital: $\rho = 1.0\times$, $0.75\times$, $0.5\times$ of base correlations.

\section{Implementation}

\textbf{Code location:} \texttt{infrastructure/risk/regulatory\_capital.py} --- SA RWA calculator.

\textbf{Sensitivities source:} \texttt{GreeksCalculator} provides delta and vega. Curvature:
\[
  CVR_i = -V(S_i + RW_i S_i) + V(S_i - RW_i S_i) - 2V(S_i)
\]

\textbf{Capital aggregation:} Three correlation scenarios; capital $= \max$ across scenarios.

\section{Validation}

\textbf{Benchmark:} FRTB-SA capital benchmarked against ISDA SIMM for non-cleared derivatives. Material divergence ($>$20\%) triggers a model review.

\textbf{Sensitivity validation:} Delta sensitivities validated against analytic Greeks for vanilla equity options.

\section{Model Limitations}

\begin{enumerate}
  \item \textbf{Curvature proxy:} CVR computed analytically using bump-and-reprice; no full Monte Carlo revaluation.
  \item \textbf{RRAO classification:} Exotic optionality identification is manual; no automated product classifier.
  \item \textbf{DRC recovery assumptions:} Recovery rate 40\% applied uniformly; actual recovery varies by seniority.
\end{enumerate}

\section{Open Findings}

\begin{findings}
\finding{FRTB-F1}{Major}{Curvature risk aggregation not validated against independent vendor benchmark.}{Open}
\finding{FRTB-F2}{Major}{RRAO treatment for path-dependent exotics not documented; manual classification in use.}{Open}
\finding{FRTB-F3}{Minor}{DRC jump-to-default recovery rate uniformly 40\%; seniority-specific recovery not modelled.}{Open}
\end{findings}

\end{document}
